\documentclass[10pt,a4paper,dvipdfmx]{jsarticle}
\usepackage[utf8]{inputenc}
\usepackage[top=1cm, bottom=2cm, left=1.5cm, right=1.5cm]{geometry}
\usepackage[dvipdfmx]{graphicx}
\usepackage{enumitem}

% カスタムコマンド
\newcommand{\hang}[1]{\settowidth{\hangindent}{#1}#1}
\newcounter{sentnum}
\newcommand{\parallelitem}[2]{
  \refstepcounter{sentnum}
  \noindent
  \parbox[t]{0.49\textwidth}{\hang{\thesentnum.~}#1}
  \hfill
  \parbox[t]{0.49\textwidth}{\hang{\thesentnum.~}#2}
  \par
  \vspace{\baselineskip}  % ここでブロックごとの間隔を調整
}

\title{Novel Moving Steady-State Visual Evoked Potential Stimulus to\\
Assess Afferent and Efferent Dysfunction in Multiple Sclerosis\\[4pt]
}
\author{大森 俊弥}
\date{2026/07/07}

\begin{document}
\maketitle

% 第一引数に英文、第二引数に日本語訳を書く
\parallelitem{Afferent and efferent visual dysfunction are prominent features of multiple sclerosis (MS).}{求心性および遠心性の視覚機能障害は，多発性硬化症（Multiple Sclerosis; MS）の主要な特徴である．}

\parallelitem{Visual outcomes have been shown to be robust biomarkers of the overall disease state.}{これら2つの視覚機能を測定する指標は，MSの全般的な病態を頑健に反映することが示されている．}

\parallelitem{Unfortunately, precise measurement of afferent and efferent function is typically limited to tertiary care facilities, which have the equipment and analytical capacity to make these measurements, and even then, only a few centers can accurately quantify both afferent and efferent dysfunction.}{しかしながら，求心性および遠心性の機能を精密に測定できるのは，通常，そのための装置と解析能力を備えた三次医療機関に限られる．これらの施設のなかでも，求心性と遠心性の機能をいずれも正確に定量化できるのは，ごく一部にとどまる．}

\parallelitem{These measurements are currently unavailable in acute care facilities (ER, hospital floors).}{そのため，急性期医療の現場（救急外来や一般病棟）では，これらの機能を現時点で測定できない．}

\parallelitem{We aimed to develop a moving multifocal steady-state visual evoked potential (mfSSVEP) stimulus to simultaneously assess afferent and efferent dysfunction in MS for application on a mobile platform.}{そこで本論文では，求心性および遠心性の機能をモバイル環境で同時に評価することを目的とした．この評価を実現するために，多焦点定常状態視覚誘発電位（multifocal steady-state visual evoked potential; mfSSVEP）を誘発する移動視覚刺激を開発した．}

\parallelitem{The brain-computer interface (BCI) platform consists of a head-mounted virtual reality headset with electroencephalogram (EEG) and electrooculogram (EOG) sensors.}{刺激を呈示するブレイン・コンピュータ・インタフェース（Brain-Computer Interface; BCI）プラットフォームは，頭部装着型のバーチャルリアリティ（Virtual Reality; VR）ヘッドセットで構成される．ヘッドセットは，脳波（Electroencephalogram; EEG）センサと眼電図（Electrooculogram; EOG）センサを備える．}

\parallelitem{To evaluate the platform, we recruited consecutive patients who met the 2017 MS McDonald diagnostic criteria and healthy controls for a pilot cross-sectional study.}{このプラットフォームを評価するため，パイロット横断研究を実施した．この研究では，一定期間内に2017年版マクドナルド診断基準を満たしたMS患者を，選別せずに全員対象とした．比較のための健常対照者も対象とした．}

\parallelitem{Nine MS patients (mean age 32.7 years, SD 4.33) and ten healthy controls (24.9 years, SD 7.2) completed the research protocol.}{このうち，研究プロトコルを完了したのは，MS患者9名（平均年齢32.7歳，標準偏差4.33）と健常対照者10名（平均年齢24.9歳，標準偏差7.2）であった．}

\parallelitem{The afferent measures based on mfSSVEPs showed a significant difference between the groups (signal-to-noise ratio of mfSSVEPs for controls: $2.50 \pm 0.72$ vs. MS: $2.04 \pm 0.47$) after controlling for age ($p = 0.049$).}{mfSSVEPに基づく求心性機能の指標は，年齢を調整した解析において，両群間で有意な差を示した（mfSSVEPの信号対雑音比は，対照群とMS患者群でそれぞれ$2.50 \pm 0.72$，$2.04 \pm 0.47$，$p = 0.049$）．}

\parallelitem{In addition, the moving stimulus successfully induced smooth pursuit movement that can be measured by the EOG signals.}{さらに，移動刺激は，EOG信号によって測定できる滑動性追従眼球運動を誘発した．}

\parallelitem{There was a trend for worse smooth pursuit tracking in cases vs.\ controls, but this did not reach nominal statistical significance in this small pilot sample.}{この追従運動では，MS患者群は対照群に比べて，移動する標的をなめらかに追従できない傾向を示した．ただし，今回の小規模なパイロット標本では，統計的に有意な差は認められなかった．}

\parallelitem{This study introduces a novel moving mfSSVEP stimulus for a BCI platform to evaluate neurologic visual function.}{本論文は，神経学的な視覚機能を評価するBCIプラットフォーム向けに，mfSSVEPを誘発する新規の移動視覚刺激を提案した．}

\parallelitem{The moving stimulus showed a reliable capability to assess both afferent and efferent visual functions simultaneously.}{この移動刺激により，求心性と遠心性の視覚機能を，同時にかつ高い信頼性で評価できることを示した．}

\end{document}
